\section{Experiments and Results}
\label{sec:results}

\subsection{Study Cohort}
\label{sec:results-cohort}

At the two-hour landmark, \RVHtwoTestEligible{} of the \RVHtwoTestStays{} test-period stays were eligible and \RVHtwoTestPositives{} had the primary outcome (Table~\ref{tab:cohort}).
Table~\ref{tab:cohort} (added analysis) compares the eligible stays of the two periods. Age, sex, and triage acuity differed little (standardized differences
\RVCharSmdAge{}, \RVCharSmdFemale{}, and \RVCharSmdAcuity{} for acuity 1 or 2). The primary-outcome rate rose from \RVCharTrainPrimary{} in the training period to
\RVCharTestPrimary{} in the test period, the share of stays with no ED vital sign after triage fell from \RVZeroRowsTrainPct{} to \RVZeroRowsTestPct{} (standardized
difference \RVCharSmdRowsZero{}), and the share with at least one comorbidity fell from \RVCharTrainElixPct{} to \RVCharTestElixPct{} (\RVCharSmdElix{}).

The change in recording (added analysis; Supplementary Table~\ref{S-tab:recording}) was a transition between 2013 and 2016 that ended before the test period. The share of eligible stays with no ED vital sign by two hours fell from \RVRecZeroC{} in 2011 to \RVRecZeroF{} in 2014, \RVRecZeroG{} in 2015, and \RVRecZeroH{} in 2016, and the median delay from arrival to the first charted ED vital sign fell from \RVRecDelayC{} hours in 2011 to \RVRecDelayH{} hours in 2016. Over the same years the share of stays whose first ED vital-sign row repeats the triage measurement rose from \RVRecEqTriC{} to \RVRecEqTriF{}, \RVRecEqTriG{}, and \RVRecEqTriH{}, and was \RVRecEqTriTest{} in the test period.

Removing the copied row left discrimination unchanged (added analysis; Supplementary Table~\ref{S-tab:dedup}). The row was removed for \RVDdFlaggedFit{} of stays in the fitting subset and \RVDdFlaggedTest{} in the test period. After refitting, the AUROC of the trees was \RVDdTrees{} (\RVDdTreesLo{}--\RVDdTreesHi{}), a change of \RVDdTreesChange{} (\RVDdTreesChangeLo{} to \RVDdTreesChangeHi{}) from the primary analysis, and that of logistic regression changed by \RVDdLrChange{} (\RVDdLrChangeLo{} to \RVDdLrChangeHi{}). Removing the timing features changed the AUROC of the trees by \RVNtTreesChange{} (\RVNtTreesChangeLo{} to \RVNtTreesChangeHi{}) and that of logistic regression by \RVNtLrChange{} (\RVNtLrChangeLo{} to \RVNtLrChangeHi{}), and reducing the inputs to vital-sign values and the static profile changed them by \RVVoTreesChange{} and \RVVoLrChange{}. The advantage of the trees over logistic regression was \RVOrigTreesOverLr{} (\RVOrigTreesOverLrLo{} to \RVOrigTreesOverLrHi{}) in the primary analysis and between \RVNtTreesOverLr{} and \RVVoTreesOverLr{} in the three variants. Scoring the stored Hybrid GRU-D models on test stays without the copied row, with no refit, changed the AUROC of the seed mean from \RVGrudHybridStored{} to \RVGrudHybridDd{}, a change of \RVGrudHybridChange{} (\RVGrudHybridChangeLo{} to \RVGrudHybridChangeHi{}), with seed-level changes from \RVGrudHybridSeedLow{} to \RVGrudHybridSeedHigh{}; this is below the threshold of \RVGrudRule{} set beforehand, so the refit was not run. The stored Temporal-Only models, which receive vital signs and their timing only, changed by \RVGrudTempChange{} (\RVGrudTempChangeLo{} to \RVGrudTempChangeHi{}).

\subsection{Timing of the Inputs and the Shortcut Diagnostic}
\label{sec:results-audit}

We first asked whether the inputs of the end-of-stay task respect an early prediction time, and how much of its discrimination needs no physiological value. The
end-of-stay models read the last ten vital-sign rows of the stay. \RVTruncTrainN{} training-period stays (\RVTruncTrainPct{}) and \RVTruncTestN{} test-period stays
(\RVTruncTestPct{}) had more rows and lost their earliest ones. Because the latest rows are kept, they extend past any early prediction time: in the test period,
\RVAfterOneHTest{} of stays kept a row charted more than one hour after arrival, \RVAfterTwoHTest{} after two hours, \RVAfterSixHTest{} after six hours, and
\RVAfterTwelveHTest{} after twelve hours, and \RVAfterDepRowsTest{} retained rows in \RVAfterDepStaysTest{} test stays (\RVAfterDepRowsTrain{} rows in
\RVAfterDepStaysTrain{} training-period stays) were charted after ED departure (Table~\ref{tab:timing}). The comorbidity count passed the audit: its diagnoses come only
from earlier admissions (\RVElixAccepted{} stay--admission links within \RVElixWindowDays{} days, none discharged at or after arrival), and an independent recomputation
reproduced the stored flags for all \RVElixFlagsIdentical{} stays. The \RVOwnAdmissionStays{} stays whose own admission carried an inconsistent discharge time were
corrected in the landmark cohorts. Removing the count changed the AUROC of Hybrid GRU-D by \RVElixDeltaHybrid{} and of the tree ensemble by \RVElixDeltaXgb{}
(Supplementary Table~\ref{S-tab:elixhauser}).

The workflow model, a logistic regression on ED duration, the number of measurements, and the per-vital measurement counts, reached an AUROC of \RVShortAAuroc{} (95\% CI
\RVShortALo{}--\RVShortAHi{}) on the end-of-stay task, and \RVShortABAuroc{} (\RVShortABLo{}--\RVShortABHi{}) with triage-time missingness added, against
\RVEosHybridAuroc{} for Hybrid GRU-D and \RVEosXgbAuroc{} for the tree ensemble on the whole-stay vital signs (Table~\ref{tab:shortcuts}). ED duration alone gave
\RVShortDurAuroc{}. In an added analysis, arrival information alone reached \RVIncArr{} (\RVIncArrLo{}--\RVIncArrHi{}) on the end-of-stay task and the workflow quantities with arrival information \RVIncArrWork{} (\RVIncArrWorkLo{}--\RVIncArrWorkHi{}), an increase of \RVIncWorkOverArr{} (\RVIncWorkOverArrLo{} to \RVIncWorkOverArrHi{}); ED duration added \RVIncDurOverArr{} (Supplementary Table~\ref{S-tab:arrival}). Within quartiles of ED duration, Hybrid GRU-D reached \RVDurQOneHybrid{} in the shortest and \RVDurQFourHybrid{} in the longest, with event rates
of \RVDurQOnePrev{} and \RVDurQFourPrev{} (Supplementary Table~\ref{S-tab:duration}).

In an added analysis by triage acuity (Supplementary Table~\ref{S-tab:acuity-shortcut}), the workflow model reached an AUROC of \RVAcuHighWork{}
(\RVAcuHighWorkLo{}--\RVAcuHighWorkHi{}) within acuity 1 or 2 (\RVAcuHighN{} stays, \RVAcuHighPos{} with the outcome) and \RVAcuMidWork{}
(\RVAcuMidWorkLo{}--\RVAcuMidWorkHi{}) within acuity 3 (\RVAcuMidN{} stays, \RVAcuMidPos{} with the outcome), against \RVAcuHighHyb{} and \RVAcuMidHyb{} for Hybrid
GRU-D. Over all stays, acuity alone reached \RVAcuOverallAcu{} and acuity with the workflow quantities \RVAcuOverallWorkAcu{}. Acuity 4 or 5 (\RVAcuLowN{} stays) had
\RVAcuLowPos{} stay with the outcome and was not evaluated.

\input{tab_shortcuts_arrival}

\subsection{Discrimination, Calibration, and Usefulness at the Landmark}
\label{sec:results-discrimination}

\subsubsection{Model Comparison}
Landmark values are not comparable with end-of-stay values, which come from other inputs, another outcome, and a larger population. The timing-aware tree ensemble
reached an AUROC of \RVAurocXgbT{} (95\% CI \RVLoXgbT{}--\RVHiXgbT{}), the end-to-end Hybrid GRU-D \RVAurocEndToEnd{} (\RVLoEndToEnd{}--\RVHiEndToEnd{}), and logistic
regression \RVAurocLr{} (\RVLoLr{}--\RVHiLr{}) (Table~\ref{tab:discrimination}); the ranking was the same for the secondary outcome. The Temporal-Only GRU-D reached
\RVAurocTmp{}. NEWS2 computed from the last values charted by the landmark reached \RVNewsHtwo{} (\RVNewsHtwoLo{}--\RVNewsHtwoHi{}), and \RVNewsHtwoComplete{} on stays
with all five parameters recorded; at arrival it reached \RVNewsTriage{}, and \RVNewsTriageAllMissing{} test stays had none of the five parameters.

Among the pre-specified contrasts (Fig.~\ref{fig:forest}; per-seed SDs and adjusted values in Supplementary Table~\ref{S-tab:contrasts}), the primary model scored \RVCOneDelta{} (\RVCOneLo{} to \RVCOneHi{}) relative to
the tree ensemble and \RVCTwoDelta{} (\RVCTwoLo{} to \RVCTwoHi{}) relative to logistic regression; both met the significance rule. The contrast with the Temporal-Only
model, which lacks the static features, was \RVCThreeDelta{} (\RVCThreeLo{} to \RVCThreeHi{}). The contrast with the frozen-encoder variant was \RVCFourDelta{}
(\RVCFourLo{} to \RVCFourHi{}) and did not meet the all-seeds rule. At the arrival landmark the tree ensemble exceeded logistic regression by \RVCFiveDelta{}
(\RVCFiveLo{} to \RVCFiveHi{}). The Holm-adjusted values equal the smallest value that \RVBootMain{} resamples can resolve. In an added analysis by approximate
calendar year of the test period, the AUROC ranged from \RVYearAurocHybridMin{} to \RVYearAurocHybridMax{} for the primary model, from \RVYearAurocTreesMin{} to
\RVYearAurocTreesMax{} for the tree ensemble, and from \RVYearAurocLrMin{} to \RVYearAurocLrMax{} for logistic regression, and the tree ensemble was at or above the
primary model in every year (Supplementary Table~\ref{S-tab:years}).

\input{tab_discrimination_cal}

\begin{figure}[t]
\centering
\includegraphics[width=\linewidth]{fig_forest}
\caption{Pre-specified AUROC contrasts on the test period (seed-mean difference with the patient-cluster bootstrap 95\% interval). Filled markers meet the significance rule; the open marker does not. A negative difference favors the second model of the pair. MIMIC-IV-ED, retrospective single-center proxy.}
\label{fig:forest}
\end{figure}

\subsubsection{Landmark Time (Added Analysis)}
\label{sec:results-sweep}
On the eligible population of each landmark (Fig.~\ref{fig:sweep}, Supplementary Table~\ref{S-tab:sweep}), the AUROC of the trees rose from arrival to \RVSweepTreesOne{} at one hour
and \RVAurocXgbT{} at two hours, and then fell to \RVSweepTreesThree{}, \RVSweepTreesFour{}, and \RVSweepTreesSix{} at three, four, and six hours; at six hours
\RVSweepEligSix{} test stays were eligible and the event rate was \RVSweepPrevSix{}. The trees exceeded logistic regression at every landmark. On the fixed set of
\RVFixN{} test stays eligible at every landmark (Supplementary Table~\ref{S-tab:fixed-cohort}), the AUROC of the trees rose from \RVFixTreesOne{}
(\RVFixTreesOneLo{}--\RVFixTreesOneHi{}) at one hour to \RVFixTreesTwo{}, \RVFixTreesThree{}, \RVFixTreesFour{}, and \RVFixTreesSix{}
(\RVFixTreesSixLo{}--\RVFixTreesSixHi{}) at two, three, four, and six hours, and that of logistic regression from \RVFixLrOne{} to \RVFixLrSix{}, while the event rate
changed from \RVFixPrevOne{} to \RVFixPrevSix{}. The arrival-landmark models were also scored on the \RVCommonN{} test stays eligible at the two-hour landmark, against the two-hour outcome (Supplementary Table~\ref{S-tab:arrival}): the arrival trees reached \RVCommonArrTrees{} (\RVCommonArrTreesLo{}--\RVCommonArrTreesHi{}) and the two-hour trees \RVCommonTwoTrees{}, a gain of \RVCommonTreesGain{} (\RVCommonTreesGainLo{} to \RVCommonTreesGainHi{}); for logistic regression the gain was \RVCommonLrGain{}.

\begin{figure}[t]
\centering
\includegraphics[width=\linewidth]{fig_sweep}
\caption{AUROC of logistic regression and timing-aware trees at landmarks from arrival to six hours (test period, primary outcome; mean over seeds with the patient-cluster bootstrap 95\% interval). Solid lines: the eligible population of each landmark, which differs between landmarks. Dashed lines: one fixed set of stays eligible at every landmark from one to six hours. MIMIC-IV-ED, retrospective single-center proxy.}
\label{fig:sweep}
\end{figure}

\subsubsection{Calibration, Subgroups, and Decision Curve (Added Analyses)}
\label{sec:results-calibration}
The primary model had a Brier score of \RVCalBrierEndToEnd{}, an expected calibration error of \RVCalEceEndToEnd{}, and a calibration slope of \RVCalSlopeEndToEnd{}
(Supplementary Table~\ref{S-tab:calibration}). Its mean predicted risk was \RVCalMeanPredEndToEnd{} against an observed rate of \RVCalObserved{} (calibration intercept
\RVCalInterceptEndToEnd{}). The reliability curves of the three best models were similar (Fig.~\ref{fig:reliability}a): close to the diagonal in the lowest decile
(\RVReliLowPred{} predicted against \RVReliLowObs{} observed for the primary model), with predicted risk below the observed rate in the upper deciles (\RVReliNinthPred{}
against \RVReliNinthObs{} in the ninth decile, \RVReliTopPred{} against \RVReliTopObs{} in the highest).

AUROC was \RVSubFemaleAuroc{} (\RVSubFemaleLo{}--\RVSubFemaleHi{}) for women and \RVSubMaleAuroc{} (\RVSubMaleLo{}--\RVSubMaleHi{}) for men, \RVSubAgeUnderAuroc{}
(\RVSubAgeUnderLo{}--\RVSubAgeUnderHi{}) below 65 years and \RVSubAgeOverAuroc{} (\RVSubAgeOverLo{}--\RVSubAgeOverHi{}) at 65 years or older, and \RVSubAcuityHighAuroc{}
(\RVSubAcuityHighLo{}--\RVSubAcuityHighHi{}) for triage acuity 1 or 2 against \RVSubAcuityMidAuroc{} (\RVSubAcuityMidLo{}--\RVSubAcuityMidHi{}) for acuity 3
(Supplementary Table~\ref{S-tab:subgroups}). Acuity 4 or 5 (\RVSubAcuityLowN{} stays) had \RVSubAcuityLowPos{} stay with the outcome and was not evaluated.

All three models had a higher net benefit than treating every stay as positive at the thresholds examined (Fig.~\ref{fig:reliability}b). At a threshold of 0.05 the
net benefit was \RVNbHybridFive{} for the primary model, \RVNbTreesFive{} for the trees, and \RVNbLrFive{} for logistic regression, against \RVNbAllFive{} for
treating all stays; at 0.10 it was \RVNbHybridTen{}, \RVNbTreesTen{}, \RVNbLrTen{}, and \RVNbAllTen{}. At 0.10, the 95\% CI of the difference between the primary
model and the trees was \RVNbHybridTreesLoTen{} to \RVNbHybridTreesHiTen{}.

\begin{figure*}[t]
\centering
\includegraphics[width=\linewidth]{fig_reliability}
\caption{(a) Reliability of the three best models on the test period: mean predicted risk against observed event rate in ten equal-frequency bins of the seed-averaged score, with the patient-cluster bootstrap 95\% interval of the observed rate (logarithmic axes; the models are offset horizontally so that the intervals do not overlap; dotted line: perfect calibration). (b) Decision curve: net benefit against threshold probability for the same models (the score is used as the probability), for treating every stay as positive, and for treating none. MIMIC-IV-ED, retrospective single-center proxy.}
\label{fig:reliability}
\end{figure*}

\subsection{Conformal Prediction on the Later Period}
\label{sec:results-coverage}

\subsubsection{Marginal Coverage}
Static split-conformal score intervals, calibrated on \RVNFc{} final-calibration stays (\RVPosFc{} with the outcome), covered \RVIntCovStatic{} (SD \RVIntCovSdStatic{}
over seeds) of the \RVNTest{} test stays at a target of 0.90, above the validity threshold of 0.89, with a mean width of \RVIntWidStatic{} (Table~\ref{tab:intervals}).
ECI covered \RVIntCovEci{} at a width of \RVIntWidEci{}, and recalibration with a 10{,}000-arrival window \RVIntCovRecalTen{} at \RVIntWidRecalTen{}. By approximate
calendar year, static coverage was \RVYearCovFirst{} in the first block (\RVYearNFirst{} stays), \RVYearCovSecond{}, \RVYearCovThird{}, and \RVYearCovFourth{} in the
last block, which holds \RVYearNFourth{} stays. The patient-cluster bootstrap 95\% interval of the static coverage of the primary model was \RVStaticBootHybLo{} to \RVStaticBootHybHi{}, which includes the threshold of 0.89 (\RVStaticBootHybShare{} of resamples reached it; Supplementary Table~\ref{S-tab:static-cov}). In added analyses with other underlying models, static coverage was \RVOtherXgbStaticCov{} with the tree ensemble (\RVStaticBootXgbLo{} to \RVStaticBootXgbHi{}; \RVStaticBootXgbShare{} of resamples reached the threshold), which does not meet the validity threshold, and \RVOtherLrStaticCov{} with logistic regression, and ECI coverage was \RVOtherXgbEciCov{} and \RVOtherLrEciCov{} (Supplementary Table~\ref{S-tab:other-models}); after
\RVStreamPerm{} within-year permutations of the arrival order per seed, ECI coverage stayed between \RVStreamEciMin{} and \RVStreamEciMax{} and recalibration between
\RVStreamRecalMin{} and \RVStreamRecalMax{} (Supplementary Table~\ref{S-tab:stream}).

\input{tab_intervals}

\subsubsection{Coverage of the Patients Who Deteriorated}
Marginal coverage did not extend to the patients who deteriorated (Table~\ref{tab:conformal}). Class-1 coverage was \RVSetCovOneStaticLac{} for static LAC sets and
\RVSetCovOneEciSets{} for ECI sets. A stay enters the marginal set with label 1 only if its predicted risk is at least \RVScoreNeeded{} (Fig.~\ref{fig:hist}), whereas the
median predicted risk among stays with the outcome was \RVScoreMedEvents{} (90th percentile \RVScoreQNinetyEvents{}) and among stays without it \RVScoreMedNon{}. The
singleton $\{1\}$ was issued for \RVSetSingleStaticLac{} of stays, with a PPV of \RVSetPpvStaticLac{} and a sensitivity of \RVSetSensStaticLac{}, and
\RVSetAbstOneStaticLac{} of the patients with the outcome received an abstention. The singleton $\{1\}$ of a static set is issued exactly for the stays whose score exceeds one threshold, so its PPV and sensitivity are a point on the precision--recall curve of the score (Supplementary Fig.~\ref{S-fig:pr}; seed 42).

Mondrian calibration raised class-1 coverage to \RVSetCovOneStaticMond{} (static), \RVSetCovOneMondEci{} (Mondrian-ECI), and \RVSetCovOneMondRecal{} (recalibrated),
at mean set sizes of \RVSetSizeStaticMond{}, \RVSetSizeMondEci{}, and \RVSetSizeMondRecal{}. For Mondrian-ECI, the singleton $\{1\}$ was issued for
\RVSetSingleMondEci{} of stays, with a PPV of \RVSetPpvMondEci{} (\RVSetPpvStaticMond{} for static Mondrian-LAC) and a sensitivity of \RVSetSensMondEci{}, and
\RVSetFaMondEci{} stays per 100 patients received it without the outcome. These quantities were similar for the static and the online class-conditional sets. For static Mondrian-LAC the PPV was \RVPrMondPpv{} at a sensitivity of \RVPrMondSens{}, against \RVPrMondSame{} for a plain threshold at that sensitivity, and for Mondrian-ECI \RVPrEciPpv{} at \RVPrEciSens{}, against \RVPrEciSame{}.

\input{tab_conformal}

\begin{figure}[t]
\centering
\includegraphics[width=\linewidth]{fig_hist}
\caption{Predicted risk of the primary model on the test period (seed 42), for stays with and without the outcome. Dashed line: the risk a stay needs to include label 1 in a marginal LAC set; dotted lines: the thresholds of the class-conditional (Mondrian) sets for label 1 (lower) and label 0 (upper). MIMIC-IV-ED, retrospective single-center proxy.}
\label{fig:hist}
\end{figure}

\subsubsection{Coverage Levels, Label Delay, and Drift}
\label{sec:results-alpha}
Across nominal coverage levels (Fig.~\ref{fig:alpha}, Supplementary Table~\ref{S-tab:alpha}), static LAC covered \RVATwentyCovOneStaticLac{} of the patients with the outcome at 0.80
and \RVAOneCovOneStaticLac{} at 0.99. Mondrian-ECI covered \RVATwentyCovOneMondEci{} of them at 0.80, with a PPV of \RVATwentyPpvMondEci{} and \RVATwentyFaMondEci{}
false alarms per 100 patients; at 0.95 its PPV was \RVAFivePpvMondEci{} with \RVAFiveFaMondEci{} false alarms and a mean set size of \RVAFiveSizeMondEci{}, and at 0.99
its PPV was \RVAOnePpvMondEci{} with \RVAOneFaMondEci{} false alarm per 100 patients and a mean set size of \RVAOneSizeMondEci{}. The class-1 quantile was feasible at
every level.

In an added analysis, shifting the intercept of the score of the primary model on the most recent labelled arrivals raised its mean predicted risk from \RVIrMeanPredOrig{} to \RVIrMeanPredRecal{} (observed \RVIrObserved{}) and lowered the Brier score from \RVIrBrierOrig{} to \RVIrBrierRecal{} (mean over five seeds; Supplementary Table~\ref{S-tab:intercept}). When the same shift was applied to the final-calibration scores and the static quantile was recomputed, marginal coverage was \RVIrCovOne{} against \RVIrCovOrig{} without recalibration, and the class-1 coverage of marginal sets \RVIrOneOne{} against \RVIrOneOrig{}; with a quantile recomputed every 2{,}000 arrivals the values were \RVIrCovTwo{} and \RVIrOneTwo{}. Static Mondrian sets covered \RVIrMondOne{} of the patients with the outcome.

\label{sec:results-robustness}
Delaying labels by 60 to 1{,}000 arrivals left ECI coverage between \RVLagSixtyEciCov{} and \RVLagThousandEciCov{}, recalibration at \RVLagSixtyRecalCov{}, and the
class-1 coverage of Mondrian-ECI at \RVLagSixtyMondEciCovOne{} and \RVLagThousandMondEciCovOne{} for the shortest and longest delay (Table~\ref{tab:lag-drift}). Under imposed drift, static coverage fell to \RVDriftShiftOneStaticPost{} after a logit
shift of $+1.0$ and from \RVDriftCasemixStaticPre{} to \RVDriftCasemixStaticPost{} after the case-mix change. ECI held \RVDriftShiftOneEciPost{} and
\RVDriftCasemixEciPost{}, whereas recalibration held the logit shifts but reached \RVDriftCasemixRecalPost{} under the case-mix change.

\begin{figure*}[t]
\centering
\includegraphics[width=\linewidth]{fig_alpha}
\caption{Class-1 coverage, positive predictive value of the singleton set $\{1\}$, and false alarms per 100 patients against the nominal coverage $1-\alpha$, for the primary model on the test period (mean over five seeds).
Dotted line: class-1 coverage equal to the nominal level. MIMIC-IV-ED, retrospective single-center proxy.}
\label{fig:alpha}
\end{figure*}

\input{tab_lag_drift}
